% Integration fragment. The notation h_a denotes the positive-first-step density;
% g_a denotes the fair-first-step density, as in sec_bulk_applications.tex.
% The existing fixed-memory weak bridge, density/unimodality inputs, and the
% certified Landau mode are explicit prerequisites, not new asserted hypotheses.

\begin{theorem}[The two peaks near memory one]\label{thm:near-one-profile}
Let $h_a$ denote the density of $\mathcal L_a$, and put
\[
 A_\varepsilon=1+2\varepsilon\log(1/\varepsilon),\qquad
 K_\varepsilon(z)=2\varepsilon
       h_{1-2\varepsilon}(A_\varepsilon-2\varepsilon z),
\]
\[
 H_\varepsilon(z)=4\varepsilon
       g_{1-2\varepsilon}(A_\varepsilon-2\varepsilon z)
       \mathbf1\{z<A_\varepsilon/(2\varepsilon)\}.
\]
As $\varepsilon\downarrow0$, both profiles converge uniformly to $f_L$,
and $\|K_\varepsilon-f_L\|_1\to0$. If $c_a>0$ is the positive mode of $g_a$,
then
\[
 c_{1-2\varepsilon}
   =1+2\varepsilon\log(1/\varepsilon)
      -2\varepsilon\lambda_{\mathrm L}+o(\varepsilon),\qquad
 4\varepsilon g_{1-2\varepsilon}(c_{1-2\varepsilon})
      \longrightarrow f_L(\lambda_{\mathrm L}).
\]
For the initial bias $q\in[0,1]$, write
$g_{a,q}(x)=q h_a(x)+(1-q)h_a(-x)$. Then
\[
 \sup_{q\in[0,1]}\sup_{x\in\mathbb R}
 \left|2\varepsilon g_{1-2\varepsilon,q}(x)
 -q f_L\!\left(\frac{A_\varepsilon-x}{2\varepsilon}\right)
 -(1-q)f_L\!\left(\frac{A_\varepsilon+x}{2\varepsilon}\right)\right|
 \longrightarrow0.
\]
The same two-component approximation holds in density $L^1$, uniformly
over $q\in[0,1]$, after dividing the displayed approximation by
$2\varepsilon$.
\end{theorem}

\begin{proof}
Write $f=f_L$, $m=\lambda_{\mathrm L}$ and
$B_\varepsilon=A_\varepsilon/(2\varepsilon)\to\infty$.
The inputs about $f$ are that it is continuous, nonnegative, integrates to
one, tends to zero at both infinities, has a unique strict maximum at $m$,
and is strictly concave on a neighbourhood of $m$. The certified mode
lemma supplies the last two assertions; the Fourier representation and
the Riemann--Lebesgue lemma supply the tail limits.

The fixed-memory bridge gives convergence of every bounded interval mass
of $K_\varepsilon$ to the corresponding interval mass of $f$.
The density $K_\varepsilon$ is unimodal, since $h_{1-2\varepsilon}$ is
unimodal~\cite{glrs2025}. Its modes tend to $m$. To see this without assuming
pointwise convergence, fix $b<m$. Choose a short interval about $m$ and
an equal-length interval strictly between $b$ and $m$. Continuity and the
strict maximum make the limiting mass of the former larger. A mode to
the left of $b$ would make the approximating density nonincreasing on
both intervals and force the reverse mass ordering. The argument to
the right is identical. On any fixed interval avoiding $m$, the
approximating densities are consequently monotone in the correct
direction for all sufficiently small $\varepsilon$. Their values are
squeezed between adjacent interval averages. Letting the interval length
decrease to zero proves
$K_\varepsilon(z)\to f(z)$ for every $z\ne m$.
The identity
\[
 \int|K_\varepsilon-f|
   =2-2\int\min(K_\varepsilon,f)
\]
and dominated convergence now give $L^1$ convergence.

If $k$ is a nonnegative unimodal density with mode $b$, then
\[
 k(y)\le \frac1{y-b}\qquad(y>b),
\]
because its integral over $[b,y]$ is at least $(y-b)k(y)$.
Since the modes of $K_\varepsilon$ remain bounded,
$\sup_{y\ge B_\varepsilon}K_\varepsilon(y)\to0$.
The fair-density identity gives, for $z<B_\varepsilon$,
\[
 H_\varepsilon(z)
   =K_\varepsilon(z)+K_\varepsilon(2B_\varepsilon-z).
\]
It follows that $H_\varepsilon(z)\to f(z)$ for every $z\ne m$.

For small enough $\varepsilon$, the strict-shape theorem gives one
positive hump of $g_{1-2\varepsilon}$ with a concavity interval between
two convex intervals. After the affine change of variables,
$H_\varepsilon$ is unimodal and has the same convex--concave--convex
partition up to $B_\varepsilon$. We make no convexity assertion across
the truncation point. The preceding mode argument, or its pointwise
version, shows that its mode $m_\varepsilon$ tends to $m$.

Choose two fixed triples, one on each side of $m$, in the neighbourhood
where $f$ is strictly concave. Each strict midpoint-concavity inequality
persists for $H_\varepsilon$. Neither convex part of
$H_\varepsilon$ can contain its corresponding triple. Hence its
concavity interval contains a fixed neighbourhood $[L,R]$ of $m$ for
all sufficiently small $\varepsilon$.
For fixed $L<\ell<x<m$ in this neighbourhood, concavity gives
\[
 H_\varepsilon(m_\varepsilon)
 \le H_\varepsilon(x)
    +\frac{m_\varepsilon-x}{x-\ell}
       \{H_\varepsilon(x)-H_\varepsilon(\ell)\}.
\]
First let $\varepsilon\downarrow0$ and then $x\uparrow m$. This bounds
the limiting upper peak height by $f(m)$. The inequality
$H_\varepsilon(m_\varepsilon)\ge H_\varepsilon(x)$ gives the reverse
bound. Thus $H_\varepsilon(m_\varepsilon)\to f(m)$.
This secant argument is necessary: weak convergence and unimodality
alone need not bound the height at the mode.

Monotonicity and a finite grid give uniform convergence on compact
intervals away from $m$. Near $m$, the converging peak height and the
values at two nearby endpoints give the same conclusion. Beyond a
large fixed interval, monotonicity bounds $H_\varepsilon$ by its endpoint
values, and $f$ tends to zero. Thus $H_\varepsilon\to f$ uniformly on
$\mathbb R$. The displayed folding identity and the uniform tail bound
then give $K_\varepsilon\to f$ uniformly below $B_\varepsilon$; above
$B_\varepsilon$ both functions vanish uniformly in the limit.

Finally,
$m_\varepsilon=(A_\varepsilon-c_{1-2\varepsilon})/(2\varepsilon)$ and
$H_\varepsilon(m_\varepsilon)
 =4\varepsilon g_{1-2\varepsilon}(c_{1-2\varepsilon})$.
These identities give the stated mode and height limits.
Reflection and mixing do not increase either the supremum error or
the density $L^1$ error: for $0\le q\le1$,
\[
 \big\|q(K_\varepsilon-f)
 +(1-q)(K_\varepsilon-f)(2B_\varepsilon-\cdot)\big\|_r
 \le\|K_\varepsilon-f\|_r,\qquad r=1,\infty.
\]
An affine change of variables gives the two assertions uniform in $q$.
\end{proof}

% Exact two-mode claims for biased starts are intentionally not in this
% fragment until the finite-window and no-three-local-maxima proofs are joined.
